\documentclass{article}
\usepackage{amsmath}
\usepackage{amssymb}
\usepackage{amsthm}
\usepackage{mathtools}
\usepackage{geometry}
\usepackage{hyperref}

\newcommand{\rk}{\operatorname{rank}}
\newcommand{\GL}{\operatorname{GL}}
\newcommand{\R}{\mathbb{R}}

\theoremstyle{definition}
\newtheorem{theorem}{Theorem}[section]
\newtheorem{proposition}[theorem]{Proposition}
\newtheorem{lemma}[theorem]{Lemma}
\newtheorem{corollary}[theorem]{Corollary}
\newtheorem{definition}[theorem]{Definition}
\newtheorem{remark}[theorem]{Remark}

\title{Example Checks for the Indecomposables}
\date{}

\begin{document}
\maketitle

\section{Goal}

This file checks the lowest nontrivial cyclic example used in the second
conjecture block:

\begin{itemize}
\item depth $L=3$, so the quiver is the oriented $4$-cycle;
\item the ambient free-arrow algebra $A_3^{\mathrm{wt}}$;
\item the fixed-arrow specialization relevant for a full-rank DLN, i.e.
  $r=d_L$ and therefore $W_0=0$.
\end{itemize}

The point of the check is to verify concretely that:

\begin{itemize}
\item the free-arrow indecomposables are exactly the admissible interval modules
  predicted by the general theory;
\item the fixed-arrow constraint is exactly the scalar condition
  $\rk(m)=\rk(W_0)$;
\item in the full-rank case $W_0=0$, the wrapped indecomposables disappear and
  one recovers the open-chain type-$A$ family.
\end{itemize}

\section{The depth-$3$ weight-space algebra}

For $L=3$, the cyclic quiver is
\[
0 \xrightarrow{a_1} 1 \xrightarrow{a_2} 2 \xrightarrow{a_3} 3 \xrightarrow{a_0} 0.
\]
The weight-space ideal is
\[
I_3^{\mathrm{wt}}
=
\langle p_1,p_2,p_3\rangle,
\]
with
\[
p_1=a_0a_3a_2,
\qquad
p_2=a_1a_0a_3,
\qquad
p_3=a_2a_1a_0.
\]
The surviving open-chain path is
\[
p_0=a_3a_2a_1.
\]
Thus
\[
A_3^{\mathrm{wt}}=kC_4/I_3^{\mathrm{wt}}.
\]

\section{Free-arrow indecomposables}

\begin{proposition}[Admissible pairs for $L=3$]
For $L=3$, the admissible index set is
\[
\mathcal I_3^{\mathrm{wt}}
=
\{(i,j)\mid 0\le i\le j\le 3\}
\cup
\{(i,j)\mid 1\le i\le 3,\ 0\le j\le i-2\}.
\]
Explicitly, this is the set of $13$ pairs
\[
\begin{aligned}
&(0,0),(0,1),(0,2),(0,3),\\
&(1,1),(1,2),(1,3),\\
&(2,2),(2,3),(2,0),\\
&(3,3),(3,0),(3,1).
\end{aligned}
\]
The forbidden wrapped endpoints are
\[
(1,0),\qquad (2,1),\qquad (3,2).
\]
\end{proposition}

\begin{proof}
The general formula gives the non-wrapping family
\[
0\le i\le j\le 3,
\]
which contributes $10$ pairs, and the wrapped family
\[
1\le i\le 3,\qquad 0\le j\le i-2,
\]
which contributes the three additional pairs
\[
(2,0),\qquad (3,0),\qquad (3,1).
\]
The forbidden pairs are exactly the excluded full wrapped intervals
\[
(i,i-1),
\qquad
1\le i\le 3.
\]
\end{proof}

\begin{corollary}[Free-arrow indecomposable list]
The indecomposable $A_3^{\mathrm{wt}}$-modules are exactly the interval modules
\[
M(i,j)
\qquad
((i,j)\in \mathcal I_3^{\mathrm{wt}}),
\]
namely:

\begin{itemize}
\item the $10$ non-wrapping modules
  \[
  M(0,0),M(0,1),M(0,2),M(0,3),
  \]
  \[
  M(1,1),M(1,2),M(1,3),
  \]
  \[
  M(2,2),M(2,3),
  \]
  \[
  M(3,3);
  \]
\item the $3$ wrapped modules
  \[
  M(2,0),\qquad M(3,0),\qquad M(3,1).
  \]
\end{itemize}
\end{corollary}

\section{The closing-arrow rank formula for $L=3$}

\begin{definition}[Closing-arrow activity for $L=3$]
Define
\[
\epsilon(i,j)=1
\]
if the closing arrow $a_0:3\to 0$ acts nontrivially on $M(i,j)$, and
\[
\epsilon(i,j)=0
\]
otherwise.
\end{definition}

\begin{proposition}[Explicit activity table]
For $L=3$, one has
\[
\epsilon(i,j)=1
\quad\Longleftrightarrow\quad
(i,j)\in \{(2,0),(3,0),(3,1)\}.
\]
Equivalently, the only closing-arrow-active indecomposables are the three
wrapped modules
\[
M(2,0),\qquad M(3,0),\qquad M(3,1).
\]
\end{proposition}

\begin{proof}
The arrow $a_0:3\to 0$ appears exactly in the wrapped intervals that pass
through the edge $3\to 0$. For $L=3$, these are precisely
\[
[2\to 0],\qquad [3\to 0],\qquad [3\to 1].
\]
No non-wrapping interval uses the closing arrow.
\end{proof}

\begin{corollary}[Rank formula for multiplicity patterns when $L=3$]
If
\[
m=(m(i,j))\in \mathcal M_d^+(A_3^{\mathrm{wt}}),
\]
then
\[
\rk(m)=m(2,0)+m(3,0)+m(3,1).
\]
\end{corollary}

\begin{proof}
This is the general formula
\[
\rk(m)=\sum \epsilon(i,j)\,m(i,j)
\]
with the explicit activity table above.
\end{proof}

\section{Full-rank DLN case: fixed arrow $W_0=0$}

\begin{proposition}[Full-rank specialization]
Assume $r=d_L$. Then
\[
P_Q=0,
\qquad
W_0=\Sigma_{XY}P_Q=0.
\]
Hence the fixed-arrow admissibility condition becomes
\[
\rk(m)=0.
\]
\end{proposition}

\begin{proof}
If $r=d_L$, then the complementary projector $P_Q$ is zero, so the closing map
vanishes. The fixed-arrow classification therefore requires
\[
\rk(m)=\rk(W_0)=0.
\]
\end{proof}

\begin{corollary}[Fixed-arrow indecomposables for $L=3$ and $r=d_L$]
In the full-rank DLN case, the only indecomposables that occur in the fixed-arrow
slice are the $10$ non-wrapping modules
\[
M(0,0),M(0,1),M(0,2),M(0,3),
\]
\[
M(1,1),M(1,2),M(1,3),
\]
\[
M(2,2),M(2,3),
\]
\[
M(3,3).
\]
The three wrapped modules
\[
M(2,0),\qquad M(3,0),\qquad M(3,1)
\]
do not appear in the fixed-arrow slice.
\end{corollary}

\begin{proof}
By the rank formula,
\[
\rk(m)=m(2,0)+m(3,0)+m(3,1).
\]
The condition $\rk(m)=0$ therefore forces
\[
m(2,0)=m(3,0)=m(3,1)=0.
\]
Hence every fixed-arrow module is a direct sum of the non-wrapping
indecomposables only.
\end{proof}

\begin{remark}[Consistency with the DLN picture]
When $r=d_L$, the cyclic equations become
\[
W_{<\ell}W_0W_{>\ell}=0
\]
with
\[
W_0=0,
\]
so they are automatic and
\[
\Gamma_{d,S}=\Omega.
\]
On the representation-theoretic side, setting the closing arrow to zero removes
all indecomposables on which $a_0$ acts nontrivially. The remaining
indecomposables are exactly the interval modules of the open chain
\[
0\to 1\to 2\to 3.
\]
Thus the fixed-arrow classification agrees with the expected full-rank limit.
\end{remark}

\section{Deliverables}

All intended checks pass in this example:

\begin{itemize}
\item checked: for $L=3$, the free-arrow algebra has exactly $13$ indecomposable
  interval modules;
\item checked: the closing-arrow-active family is exactly
  \[
  \{M(2,0),M(3,0),M(3,1)\};
  \]
\item checked: for $L=3$, the scalar rank formula is
  \[
  \rk(m)=m(2,0)+m(3,0)+m(3,1);
  \]
\item checked: in the full-rank DLN case $r=d_L$, the fixed-arrow condition
  $\rk(m)=0$ removes all wrapped indecomposables and leaves exactly the
  type-$A_4$ family.
\end{itemize}

Therefore the depth-$3$ sanity check agrees with the general formulas used in
the second-conjecture block.

\end{document}
